% ═══════════════════════════════════════════════════════════════════
%  Chapitre 4 — Mise en œuvre
% ═══════════════════════════════════════════════════════════════════
\chapter{Mise en \oe uvre}

\section*{Introduction}
\addcontentsline{toc}{section}{Introduction}

Ce chapitre présente la réalisation effective de la solution. Il décrit d'abord
la structure du code des trois applications, puis la stratégie de tests mise en
place et son automatisation par intégration continue. Il présente ensuite les
interfaces développées, illustrées par des captures d'écran, et se termine par
les modalités de déploiement et d'exploitation de la plate-forme.

% ─────────────────────────────────────────────────────────────────
\section{Structure du projet}

Le projet est organisé en mono-dépôt : les trois applications vivent dans un
même dépôt Git, à trois emplacements distincts. Ce choix évite la
désynchronisation entre le contrat de l'API et ses consommateurs --- une
modification du serveur et son adaptation côté client sont validées ensemble
par la même chaîne d'intégration continue.

\begin{table}[H]
\centering
\caption{Organisation du mono-dépôt}
\label{tab:monodepot}
\begin{tabularx}{\textwidth}{L{2.4cm} L{3.4cm} X}
\toprule
\entete{Dossier} & \entete{Rôle} & \entete{Pile technique} \\
\midrule
\texttt{backend/} & API REST, règles métier, base de données &
NestJS · Prisma · PostgreSQL \\
\texttt{backoffice/} & Interface de gestion &
React · Vite · Zustand \\
\texttt{mobile/} & Application cliente &
Expo · React Native · i18next \\
\texttt{scripts/} & Sauvegarde de la base & Shell \\
\texttt{brand/} & Identité visuelle et génération des icônes & Python \\
\bottomrule
\end{tabularx}
\end{table}

\subsection{Le backend}

Le code de l'API représente environ 8~000 lignes, auxquelles s'ajoutent 5~150
lignes de tests. Il est découpé en quatorze modules NestJS, réunis par un
module racine.

\figProjet[0.7]{structure-backend.png}{Structure des modules du backend}{structure-backend}

Douze d'entre eux exposent un contrôleur et suivent tous la même organisation,
ce qui rend le code prévisible :

\begin{itemize}[leftmargin=1.2cm]
  \item le \textbf{contrôleur} déclare les routes, leurs droits d'accès, leurs
  limites de débit et leur documentation ; il ne contient aucune règle ;
  \item le \textbf{service} porte la logique métier et les accès aux données ;
  \item les \textbf{objets de transfert} décrivent et valident les entrées ;
  \item les \textbf{tests} accompagnent le module qu'ils vérifient.
\end{itemize}

Les deux modules restants n'exposent aucune route : \texttt{prisma} met le
client de base de données à la disposition des autres, et \texttt{mail}
l'acheminement des courriels.

Trois dossiers, enfin, ne sont pas des modules mais du code partagé :
\texttt{common} regroupe le filtre d'exceptions, l'intercepteur de traduction,
la pagination et les validateurs communs, \texttt{config} valide les variables
d'environnement au démarrage, et \texttt{i18n} porte les messages traduits de
l'API.

\paragraph{Le c\oe ur du système.} Trois services concentrent l'essentiel de la
complexité : le service des rendez-vous (693 lignes), qui porte le moteur de
créneaux, la gestion du fuseau et le verrou de réservation ; le service de
fidélité (837 lignes), qui porte les deux circuits de récompense et
l'idempotence des crédits ; et le service d'authentification (477 lignes), qui
porte la rotation des jetons et les défenses contre l'énumération de comptes.

\subsection{Le back-office}

Le back-office représente environ 6~500 lignes de TypeScript et React,
organisées en huit pages et sept composants partagés.

\figProjet[0.7]{structure-backoffice.png}{Structure du back-office}{structure-backoffice}

\begin{itemize}[leftmargin=1.2cm]
  \item \texttt{pages/} --- une page par domaine métier : tableau de bord,
  commandes, rendez-vous, catalogue, fidélité, horaires, utilisateurs,
  connexion ;
  \item \texttt{components/} --- les éléments réutilisés : mise en page,
  pagination, boîte de confirmation, garde de route authentifiée, formulaire de
  catalogue, fenêtre d'export, fenêtre de réservation pour une cliente ;
  \item \texttt{api/} --- une façade par domaine, seule à connaître les
  adresses des routes ; les pages ne manipulent jamais d'URL ;
  \item \texttt{stores/} --- l'état global, réduit à la session authentifiée ;
  \item \texttt{theme/} --- la charte graphique partagée.
\end{itemize}

Le client HTTP centralise deux comportements délicats : l'ajout automatique du
jeton d'accès à chaque requête, et le rafraîchissement transparent de la
session lorsqu'il expire --- la requête ayant échoué est alors rejouée, sans
que l'utilisatrice s'aperçoive de quoi que ce soit.

\subsection{L'application mobile}

L'application mobile représente environ 9~700 lignes, réparties en vingt-quatre
écrans et dix-huit composants.

\figProjet[0.7]{structure-mobile.png}{Structure de l'application mobile}{structure-mobile}

La navigation s'organise en une barre d'onglets --- Accueil, Prestations,
Produits, Fidélité, Profil --- surmontant des piles de navigation propres à
chaque section. Le panier et la session sont conservés dans des magasins
d'état, le jeton de session étant stocké dans le coffre sécurisé du système
d'exploitation plutôt qu'en clair.

L'application est traduite en trois langues, avec chargement de la langue du
téléphone au démarrage. La prise en charge de l'arabe a demandé une attention
particulière : au-delà de la traduction, elle impose de vérifier chaque écran
dans le sens d'écriture de droite à gauche.

% ─────────────────────────────────────────────────────────────────
\section{Tests et validation}

\subsection{Stratégie de tests}

La stratégie retenue repose sur un constat simple : \emph{les propriétés les
plus importantes de ce projet ne sont pas démontrables par un test unitaire
classique}. Qu'un calcul de points donne le bon résultat se vérifie avec une
base simulée. Que trois réservations simultanées ne puissent pas occuper deux
cabines ne se vérifie qu'avec une vraie base, un vrai verrou et de vraies
transactions.

Les tests se répartissent donc en deux familles.

\begin{itemize}[leftmargin=1.2cm]
  \item \textbf{Les tests unitaires}, sur base simulée. Ils vérifient les
  règles métier pures : calcul des points, transitions d'états autorisées,
  génération de la grille de créneaux, validation des entrées, contrôle des
  droits.

  \item \textbf{Les tests d'intégration}, sur une véritable base PostgreSQL.
  Ils vérifient ce qu'une base simulée ne saurait démontrer : sérialisation des
  réservations concurrentes, atomicité du crédit de fidélité, annulation
  effective d'une transaction en échec, idempotence de l'émission des bons,
  rotation et détection de réutilisation des jetons, anonymisation des comptes.
\end{itemize}

Au total, le backend est couvert par 231 tests répartis en 25 suites, dont 8
écrivent réellement en base.

\begin{table}[H]
\centering
\caption{Répartition des suites de tests par domaine}
\label{tab:suites-tests}
\begin{tabularx}{\textwidth}{L{3.0cm} C{1.3cm} C{1.9cm} X}
\toprule
\entete{Domaine} & \entete{Suites} & \entete{Base réelle} & \entete{Propriétés vérifiées} \\
\midrule
Rendez-vous & 5 & 1 &
Grille de créneaux, fuseau du centre, fermetures, machine à états, capacité
sous concurrence \\

Fidélité & 4 & 2 &
Calcul et idempotence des crédits, débit conditionnel du solde, paliers de
visites, émission et remise des bons \\

Commandes & 5 & 2 &
Contrôle et décrément atomique du stock, restitution à l'annulation, crédit à
la remise, machine à états \\

Authentification & 3 & 1 &
Hachage, rotation des jetons, détection de réutilisation, validation du jeton
d'accès \\

Utilisateurs & 3 & 1 &
Profil, unicité de l'administratrice, anonymisation du compte supprimé \\

Exports & 2 & 0 &
Génération des classeurs, restriction à la gérante, bornes de dates \\

Horaires & 1 & 0 & Refus des plages qui se chevauchent \\
Infrastructure & 2 & 1 & Connexion à la base, sonde de santé \\
\midrule
\textbf{Total} & \textbf{25} & \textbf{8} & \textbf{231 tests} \\
\bottomrule
\end{tabularx}
\end{table}

\subsection{Le test de concurrence, en détail}

Le test le plus instructif du projet mérite d'être décrit, car il a
directement dicté une décision d'architecture.

Il place le centre à une capacité de deux places, puis lance trois réservations
\emph{réellement simultanées} sur le même créneau. La propriété attendue est
que deux réussissent et qu'une échoue --- jamais trois.

\begin{lstlisting}[caption={Principe du test de concurrence sur les réservations},captionpos=b]
// Trois réservations lancées en parallèle sur le même créneau,
// avec une capacité de 2 places.
const resultats = await Promise.allSettled([
  service.create(cliente1.id, { serviceId, startAt }),
  service.create(cliente2.id, { serviceId, startAt }),
  service.create(cliente3.id, { serviceId, startAt }),
]);

const reussies = resultats.filter((r) => r.status === 'fulfilled');
expect(reussies).toHaveLength(2);   // et jamais 3
\end{lstlisting}

Sur l'implémentation initiale --- compter les rendez-vous du créneau, puis
insérer --- ce test échouait de façon reproductible : les trois réservations
étaient acceptées. Les trois comptaient avant qu'aucune n'ait inséré, et les
trois trouvaient de la place. C'est ce test, et non une revue de code, qui a
imposé le recours au verrou consultatif décrit au chapitre~2.

Ces tests partagent la même base et le même verrou : ils doivent donc être
exécutés en série, ce qui est imposé par une option dédiée du lanceur de tests.

\subsection{Intégration continue}

L'ensemble est rejoué automatiquement à chaque modification par une chaîne
GitHub Actions. Elle comporte trois travaux.

\begin{table}[H]
\centering
\caption{Travaux de la chaîne d'intégration continue}
\label{tab:ci}
\begin{tabularx}{\textwidth}{L{2.8cm} X}
\toprule
\entete{Travail} & \entete{Étapes} \\
\midrule
Backend &
Installation des dépendances, génération du client Prisma, application des
migrations sur une base PostgreSQL éphémère, exécution du script d'amorçage,
compilation, puis exécution des 25 suites de tests. \\

Back-office & Installation des dépendances et vérification de types. \\

Mobile & Installation des dépendances et vérification de types. \\
\bottomrule
\end{tabularx}
\end{table}

Un point de cette chaîne mérite d'être signalé, car il n'est pas évident : le
script d'amorçage y est indispensable. Il crée les horaires d'ouverture, et
sans une seule ligne dans cette table, le centre est considéré fermé tous les
jours --- toute la suite de tests des rendez-vous échouerait pour une raison
qui n'a rien à voir avec le code testé.

\figProjet[0.85]{ci-github.png}{Exécution de la chaîne d'intégration continue}{ci-github}

\subsection{Résultats}

\figProjet[0.85]{resultats-tests.png}{Résultat de l'exécution complète de la suite de tests}{resultats-tests}

L'exécution complète est verte : 25 suites, 231 tests, aucun échec. Le
back-office et l'application mobile, qui n'ont pas de tests automatisés, sont
couverts par la vérification de types --- un filet plus mince, mais qui détecte
effectivement toute divergence entre le contrat de l'API et ses consommateurs.

% ─────────────────────────────────────────────────────────────────
\section{Interfaces graphiques}

\subsection{Application mobile}

\subsubsection{Authentification}

\figDouble{mobile-connexion.png}{Écran de connexion}{mobile-inscription.png}{Écran d'inscription}{mobile-auth}

L'inscription recueille l'adresse, le mot de passe, le nom, le prénom et le
téléphone. Ces trois derniers champs sont obligatoires : un rendez-vous ou une
commande sans nom ni numéro de téléphone est inexploitable pour l'institut, qui
doit pouvoir rappeler la cliente. Un code de vérification est envoyé par
e-mail à l'issue de l'inscription.

\figDouble{mobile-verification.png}{Vérification de l'adresse}{mobile-mdp-oublie.png}{Mot de passe oublié}{mobile-verif}

\subsubsection{Accueil et catalogues}

\figDouble{mobile-accueil.png}{Écran d'accueil}{mobile-prestations.png}{Catalogue des prestations}{mobile-accueil-prestations}

L'accueil réunit en une vue le solde de fidélité, une sélection de prestations
et une sélection de produits. C'est l'écran le plus consulté, et celui dont le
coût en requêtes a été le plus surveillé : il ne demande que ce qu'il affiche.

\figDouble{mobile-detail-prestation.png}{Détail d'une prestation}{mobile-produits.png}{Catalogue des produits}{mobile-details}

\subsubsection{Réservation}

\figDouble{mobile-reservation.png}{Choix de la date et du créneau}{mobile-recap-reservation.png}{Récapitulatif avant validation}{mobile-reservation}

L'écran de réservation affiche la grille des créneaux du jour choisi. Les
créneaux indisponibles --- capacité atteinte, ou heure déjà passée --- restent
visibles mais désactivés, plutôt que masqués : une cliente qui ne voit pas le
créneau de 15\,h ignore s'il est pris ou si l'institut n'ouvre pas à cette
heure-là.

Un jour fermé, qu'il le soit par les horaires hebdomadaires ou par une
fermeture exceptionnelle, est annoncé explicitement.

\figDouble{mobile-confirmation-rdv.png}{Confirmation du rendez-vous}{mobile-mes-rdv.png}{Mes rendez-vous}{mobile-rdv}

\subsubsection{Commande de produits}

\figDouble{mobile-detail-produit.png}{Détail d'un produit}{mobile-panier.png}{Panier}{mobile-produit-panier}

\figDouble{mobile-checkout.png}{Validation de commande}{mobile-mes-commandes.png}{Mes commandes}{mobile-commande}

La validation propose le retrait à l'institut ou la livraison à domicile,
laquelle exige une adresse. Le paiement s'effectue à la remise. La cliente peut
annuler sa commande tant qu'elle ne lui a pas été remise.

\subsubsection{Fidélité}

\figDouble{mobile-fidelite.png}{Écran de fidélité}{mobile-recompenses.png}{Mes récompenses et bons}{mobile-fidelite}

L'écran de fidélité réunit le solde de points, la progression vers le prochain
palier de visites, le catalogue des récompenses échangeables et l'historique
des mouvements. L'écran des récompenses présente les bons dus : ceux offerts
par un palier et ceux obtenus par échange de points, avec leur code à présenter
au comptoir.

\subsubsection{Profil}

\figDouble{mobile-profil.png}{Profil et réglages}{mobile-edition-profil.png}{Modification du profil}{mobile-profil}

Les réglages permettent de changer de langue, de basculer entre les modes clair
et sombre, de modifier le mot de passe et de supprimer le compte.

\subsection{Back-office}

\figProjet{bo-connexion.png}{Back-office --- écran de connexion}{bo-connexion}

Une cliente est refusée à l'entrée du back-office : ce contrôle est effectué
par le serveur et non par l'interface.

\figProjet{bo-dashboard.png}{Back-office --- tableau de bord}{bo-dashboard}

Le tableau de bord réunit ce qui appelle une action immédiate : les commandes
en attente de confirmation, les rendez-vous du jour et les produits en stock
faible. La notion d'« aujourd'hui » y est calculée par le serveur, dans le
fuseau du centre : c'est lui qui fait autorité, et non le poste qui consulte.

\figProjet{bo-commandes.png}{Back-office --- gestion des commandes}{bo-commandes}

Les commandes sont présentées par statut. Chaque changement de statut est
soumis à la machine à états : les seules transitions proposées sont celles que
le serveur accepterait.

\figProjet{bo-rendezvous.png}{Back-office --- gestion des rendez-vous}{bo-rendezvous}

Cette page permet également de réserver pour le compte d'une cliente ---
c'est-à-dire de saisir un rendez-vous pris par téléphone --- et de reporter un
rendez-vous existant, avec les mêmes contrôles de disponibilité que ceux
appliqués à une réservation cliente.

\figProjet{bo-catalogue.png}{Back-office --- gestion du catalogue}{bo-catalogue}

Le catalogue couvre les prestations, les produits et leurs catégories, y
compris le téléversement des photos et le suivi des stocks. Un élément retiré
du catalogue est désactivé et non supprimé : il disparaît immédiatement des
applications clientes, mais reste rattaché à l'historique des commandes et des
rendez-vous qui le référencent.

\figProjet{bo-fidelite.png}{Back-office --- programme de fidélité}{bo-fidelite}

Cette page réunit le catalogue des récompenses, les paliers de visites, la
liste des bons dus au comptoir, l'ajout manuel de points --- motif obligatoire
--- et le journal d'audit de ces ajouts, consultable par la seule gérante.

\figProjet{bo-horaires.png}{Back-office --- horaires et fermetures}{bo-horaires}

Les horaires se règlent par plages, à raison de plusieurs plages par jour. Un
jour sans aucune plage est fermé. Les fermetures exceptionnelles --- congés,
jours fériés --- priment sur les horaires hebdomadaires. La capacité du centre
et l'écart entre créneaux se règlent depuis cette même page, sans
redéploiement.

\figProjet{bo-utilisateurs.png}{Back-office --- gestion des utilisateurs}{bo-utilisateurs}

\figProjet[0.7]{bo-export.png}{Back-office --- fenêtre d'export Excel}{bo-export}

Les exports produisent des classeurs Excel des clientes, des commandes et des
rendez-vous, sur une période choisie. Le fichier est destiné à la gérante et
non à un développeur : aucune valeur technique brute n'y apparaît --- les
statuts y sont écrits en toutes lettres --- et les montants y restent des
\emph{nombres} plutôt que du texte, de sorte qu'une somme ou un tableau croisé
fonctionne dans le classeur.

\subsection{Courriers électroniques}

\figProjet[0.7]{mail-verification.png}{Courriel de vérification d'adresse}{mail-verification}

Les courriels sont envoyés dans la langue choisie par la cliente. En
développement, aucun message ne part : les codes s'affichent dans les journaux
du serveur, ce qui permet de dérouler l'intégralité des parcours --- inscription,
vérification, mot de passe oublié --- sans dépendre d'un fournisseur de
messagerie.

% ─────────────────────────────────────────────────────────────────
\section{Déploiement et exploitation}

\subsection{Mise en production}

La pile complète se lance par une commande unique, qui enchaîne la base de
données, l'application des migrations et le démarrage de l'API. La
documentation interactive de l'API est automatiquement coupée en production.

\subsection{Sauvegardes}

Un script produit un cliché horodaté de la base et supprime ceux de plus de
quatorze jours, à automatiser par une tâche planifiée. Les deux règles
d'exploitation qui l'accompagnent --- recopier le cliché hors de la machine, et
copier séparément le dossier des images, qui vit hors de la base --- ont été
exposées au chapitre~4 : elles relèvent de la disponibilité des données autant
que de l'exploitation courante.

La procédure de restauration, et la vérification qui doit la suivre, sont
documentées dans le script lui-même --- à l'endroit où on la cherchera le jour
où elle servira.

\subsection{Ce qui reste à faire avant une ouverture publique}

Par honnêteté sur l'état de la livraison, trois points restent ouverts :

\begin{itemize}[leftmargin=1.2cm]
  \item l'acquisition d'un nom de domaine et d'un certificat HTTPS ;
  \item la mise en place d'un serveur de messagerie définitif, à la place de la
  solution provisoire actuelle ;
  \item les notifications \textit{push}, qui dépendent d'une compilation
  native, donc du nom de domaine.
\end{itemize}

\section*{Conclusion}
\addcontentsline{toc}{section}{Conclusion}

Ce chapitre a présenté la mise en \oe uvre de la solution : la structure du code
des trois applications, la stratégie de tests et son automatisation, les
interfaces réalisées et les modalités d'exploitation. La démonstration du
produit final à travers ses écrans montre une plate-forme complète et
cohérente ; les 231 tests qui l'accompagnent, dont ceux qui écrivent dans une
véritable base, en établissent les propriétés les plus difficiles.
